\clearpage
\section{Tema de investigación}

El proyecto propuso un sistema híbrido para la predicción rápida de incendios forestales mediante simulación multiescala, procesamiento geoespacial, computación de alto rendimiento y escenarios precalculados. Aprovechó herramientas existentes para apoyar la toma de decisiones durante una emergencia.

En la escala forestal se empleó FARSITE, capaz de estimar la propagación del fuego a partir del punto de ignición, la topografía, el combustible y las condiciones meteorológicas. Para zonas críticas de interfaz urbano-forestal (WUI) se utilizó WFDS, que ofreció una simulación detallada de la combustión, el flujo y la transferencia de calor.

La combinación permitió utilizar un modelo rápido donde no se requería detalle y reservar el modelo costoso para sectores de riesgo. Para reducir los tiempos de respuesta, se construyó una biblioteca de escenarios variando el inicio, el viento, las condiciones ambientales, el terreno y el combustible.

Ante un incendio, el usuario seleccionó su ubicación en un mapa e ingresó las condiciones. El sistema comparó los datos, recuperó escenarios similares y mostró una estimación visual de la propagación y del impacto urbano.

El aporte de Ingeniería de Sistemas incluyó integrar FARSITE y WFDS, transformar coordenadas, geometrías y parámetros, procesar datos SIG, paralelizar simulaciones con HPC y desarrollar mecanismos de almacenamiento, búsqueda y visualización. Esto convirtió resultados técnicos complejos en información rápida, accesible y útil para orientar la respuesta de los organismos de emergencia durante emergencias forestales reales y situaciones críticas de alto riesgo.
